\chapter{Các nghiên cứu liên quan}
\label{chap:relatedwork}

\section{Tổng quan}
Chương này trình bày các nghiên cứu liên quan trên các lĩnh vực chính: ứng dụng của theo dõi chuyển động, cảm biến quán tính cho HAR, các kiến trúc AI phù hợp, phương pháp tiền xử lý và trích xuất đặc trưng, kỹ thuật tối ưu hóa triển khai biên, và các khoảng trống nghiên cứu cần giải quyết.

\section{Các lĩnh vực ứng dụng của theo dõi chuyển động}
Theo dõi chuyển động con người là công nghệ nền tảng trong nhiều lĩnh vực. Trong y tế, cảm biến đeo cho phép phân tích dáng đi, phát hiện té ngã (độ nhạy 95--98\% \cite{wan2021deep_har_review}), và giám sát phục hồi chức năng \cite{article:human_motion_tracking_survey}. Thể thao tận dụng dữ liệu quán tính và sinh cơ học để tối ưu hiệu suất và phòng ngừa chấn thương. VR/AR phụ thuộc vào theo dõi khung xương chính xác cho trải nghiệm nhập vai. Ứng dụng công nghiệp gồm giám sát an toàn lao động, đạt 94,3\% trong phát hiện tư thế nguy hiểm tại công trường \cite{wang2023har_industrial}. Nhu cầu đa dạng này tạo động lực cho các framework theo dõi chuyển động hiệu quả, có khả năng suy luận biên thời gian thực với độ chính xác 90\%+.

\section{Theo dõi chuyển động dựa trên cảm biến}
Đơn vị đo lường quán tính (IMU)---kết hợp gia tốc kế và con quay hồi chuyển---được sử dụng rộng rãi trong theo dõi chuyển động nhờ tiêu thụ điện năng thấp (2--10~mA), chi phí thấp (\$1--5/đơn vị) và phù hợp để triển khai trên thiết bị biên. Filippeschi và cộng sự \cite{filippeschi2017survey_imu} tổng hợp các thách thức chính: trôi cảm biến (0,1--1$^\circ$/s cho con quay hồi chuyển), lỗi tích lũy trong ước lượng hướng, và yêu cầu hiệu chuẩn. Phương pháp hợp nhất đa cảm biến (IMU + từ kế + cảm biến áp suất) cải thiện độ chính xác 2--5\% nhưng tăng độ phức tạp và chi phí năng lượng. Xia và cộng sự \cite{xia2020lstm_har} đã chứng minh hệ thống IMU đơn đạt 96,8\% trên bài toán HAR 12 lớp với kiến trúc LSTM-CNN, xác nhận tính khả thi của cách tiếp cận cảm biến đơn.

Các tập dữ liệu chuẩn UCI-HAR \cite{anguita2013har} và WISDM \cite{banos2014wisdm} đều sử dụng IMU sáu trục gắn trên cơ thể ở 50--100~Hz, xác lập cấu hình này làm thiết lập tham chiếu cho nghiên cứu HAR trên thiết bị đeo. Framework trong luận văn kế thừa cấu hình đó---thu thập dữ liệu gia tốc kế và con quay hồi chuyển ba trục ở 100~Hz---tập trung vào quy trình IMU đơn với khả năng mở rộng đa cảm biến.

\section{Theo dõi chuyển động dựa trên thị giác}
Các phương pháp thị giác máy tính sử dụng camera và ước lượng tư thế (MediaPipe \cite{web:Google_Mediapipe}, OpenPose) cho phép tái tạo toàn thân (17+ điểm khóa, 30 FPS) và được dùng trong hoạt hình, phân tích lâm sàng. Tuy nhiên, hệ thống dựa trên thị giác gặp các hạn chế: che khuất trong cảnh đông đúc, nhạy với điều kiện ánh sáng, lo ngại quyền riêng tư khi thu hình ảnh, và tiêu thụ điện năng cao (500--2000~mW so với 20--50~mW của IMU). Đối với giám sát liên tục trên thiết bị đeo, nơi ưu tiên tuổi thọ pin (ngày đến tuần), các giải pháp IMU thuần túy vẫn là lựa chọn phù hợp nhất. Do đó, luận văn này tập trung hoàn toàn vào cảm biến quán tính.

\section{AI trong theo dõi chuyển động}
Các mô hình AI chuyển đổi dữ liệu chuyển động thô thành các nhận dạng ngữ nghĩa---phân loại hoạt động, nhận dạng tư thế và dự báo quỹ đạo. Phương pháp học máy truyền thống gồm Support Vector Machines (SVM) \cite{cortes1995svm} và Random Forests \cite{breiman2001random} đạt 88--93\% trên các benchmark chuẩn (UCI-HAR, WISDM) với bộ nhớ thời gian chạy 50--200~KB, phù hợp thiết bị biên \cite{wang2024har_comparative}.

Các kiến trúc học sâu cải thiện hiệu suất: CNN đạt 92--95\% (trích xuất đặc trưng không gian từ cửa sổ cảm biến), RNN/LSTM đạt 94--97\% (nắm bắt phụ thuộc thời gian \cite{xia2020lstm_har}), và transformer đạt 96--98\% trên tập dữ liệu phức tạp \cite{singh2023har_transformers}. Tuy nhiên, mô hình học sâu đòi hỏi tập dữ liệu lớn hơn (10.000+ mẫu so với 1.000--5.000 cho ML cổ điển), bộ nhớ cao hơn (500~KB--2~MB) và chiến lược lượng tử hóa phức tạp hơn khi triển khai biên \cite{lin2022mcunet}.

Nghiên cứu so sánh \cite{wang2024har_comparative} cho thấy ML cổ điển đạt 90--93\% với kích thước nhỏ hơn 10$\times$ và suy luận nhanh hơn 5--20$\times$ (12--18~ms so với 60--120~ms) trên Cortex-M4. Framework đề xuất hỗ trợ cả hai chiến lược---mô hình cổ điển cho kịch bản giới hạn tài nguyên, mạng nơ-ron nhẹ cho ứng dụng ưu tiên độ chính xác---với các chế độ tối ưu hóa giúp cân nhắc đánh đổi một cách có căn cứ.

\section{Các framework và nền tảng hiện có}
Như đã tổng quan trong Chương~\ref{chap:introduction}, nhiều nền tảng đã xuất hiện phục vụ phát triển mô hình AI cho thiết bị biên. Về mặt kỹ thuật, các hạn chế chính bao gồm:
\begin{itemize}
    \item \textbf{Edge Impulse} \cite{web:Edge_Impulse}: Quy trình tự động hóa cao nhưng cố định cửa sổ tiền xử lý sau khi tải dữ liệu; tinh chỉnh đòi hỏi tải lại toàn bộ. Logic tiền xử lý không minh bạch.
    \item \textbf{SensiML} \cite{web:SensiML}: AutoML với 200+ đặc trưng ứng viên nhưng sử dụng định dạng độc quyền và lựa chọn đặc trưng kiểu hộp đen.
    \item \textbf{TensorFlow Lite Micro} \cite{david2020tensorflow}: Runtime hiệu quả cho mạng nơ-ron lượng tử hóa trên vi điều khiển, nhưng chỉ hỗ trợ mạng nơ-ron và yêu cầu thiết kế kiến trúc thủ công.
    \item \textbf{Teachable Machine} \cite{web:Google_TeachableMachine}: Giao diện không mã cho tạo mẫu nhanh, nhưng thiếu trích xuất đặc trưng nâng cao và tối ưu hóa theo phần cứng.
    \item \textbf{AutoML/NAS} \cite{cai2018proxylessnas}: Tìm kiến trúc tối ưu dưới ràng buộc tài nguyên, nhưng đòi hỏi tài nguyên tính toán lớn và thiếu minh bạch.
\end{itemize}
Framework đề xuất giải quyết các hạn chế này bằng cách cung cấp: (1) tiền xử lý tương tác với phân đoạn kéo thả, (2) trích xuất đặc trưng minh bạch với kiểm tra từng đặc trưng, (3) tối ưu hóa đa mục tiêu theo ràng buộc phần cứng, và (4) mã nguồn mở không rào cản chi phí.

\section{Tiền xử lý dữ liệu và tạo cửa sổ}
Chất lượng nhận dạng hoạt động phụ thuộc nhiều vào chiến lược phân đoạn và tạo cửa sổ. Banos và cộng sự \cite{banos2014wisdm} đánh giá ảnh hưởng của kích thước cửa sổ: 1~giây đạt 89,3\%, 2,5~giây đạt 92,1\%, 5~giây đạt 93,4\% nhưng giảm khả năng phản hồi. Chồng lấp 50\% cải thiện 2--4\% so với cửa sổ không chồng lấp \cite{wan2021deep_har_review} với chi phí tính toán gấp đôi.

Lựa chọn bộ lọc ảnh hưởng đáng kể đến giảm nhiễu: bộ lọc thông thấp Butterworth (cutoff 20~Hz) bảo tồn đặc tính chuyển động trong khi làm suy giảm nhiễu cảm biến, cải thiện 3--7\% \cite{filippeschi2017survey_imu}. Tuy nhiên, quy trình tiền xử lý thông thường đòi hỏi sửa mã thủ công để điều chỉnh tham số lọc---các nhà phát triển thường dành 40--60\% thời gian dự án cho bước này \cite{merenda2020edge}. Nghiên cứu về học máy tương tác \cite{amershi2014power} chỉ ra rằng giao diện trực quan giảm 35\% thời gian lặp lại và 28\% lỗi so với quy trình dựa trên mã. Framework đề xuất áp dụng nguyên tắc này: cho phép căn chỉnh trực quan tín hiệu thô và đã xử lý, điều chỉnh cửa sổ tương tác với phản hồi tức thì, và gán nhãn nhanh.

\section{Trích xuất và biểu diễn đặc trưng}
Đặc trưng thủ công miền thời gian (trung bình, phương sai, RMS, độ lệch, độ nhọn, tỷ lệ qua không) và miền tần số (năng lượng phổ, tần số thống trị, entropy phổ) vẫn hiệu quả cho ML cổ điển khi được trích xuất có hệ thống \cite{oppenheim1999discrete}. Bộ đặc trưng điển hình gồm 60--150 đặc trưng mỗi cửa sổ \cite{wang2024har_comparative}. Lựa chọn đặc trưng (thông tin hỗ tương, kiểm định F ANOVA, loại bỏ đệ quy) có thể giảm 30--50\% chiều với mất độ chính xác dưới 2\%, đồng thời cải thiện tốc độ suy luận 15--30\% trên vi điều khiển \cite{banos2014wisdm}.

Đặc trưng miền tần số đặt ra thách thức triển khai: FFT trên hệ thống nhúng không có bộ tăng tốc phần cứng đòi hỏi O(N log N) phép toán dấu phẩy động, tiêu thụ 50--150~ms mỗi cửa sổ trên Cortex-M4 64~MHz \cite{banbury2021mlperf}. Các giải pháp nhẹ bao gồm đếm qua không, đỉnh tự tương quan, hoặc thư viện FFT tối ưu (ARM CMSIS-DSP). Quan trọng hơn, tính tương đồng đặc trưng giữa huấn luyện (Python NumPy FFT) và triển khai (C++ FFT điểm cố định) cần được xác minh cẩn thận.

Mặc dù mạng sâu end-to-end có thể học đặc trưng ngầm định, trích xuất đặc trưng tường minh cải thiện khả năng diễn giải và minh bạch triển khai---đặc biệt quan trọng trong các lĩnh vực y tế hoặc an toàn. Framework kết hợp các mô-đun đặc trưng hoán đổi được, xác minh tính nhất quán huấn luyện-triển khai, và hỗ trợ chẩn đoán từng cửa sổ.

\section{Triển khai biên và tối ưu hóa}
Ràng buộc tài nguyên định hình các thách thức triển khai biên: RAM hạn chế (32--256~KB), flash (256~KB--1~MB), không có hoặc hạn chế FPU, và ngân sách điện năng 10--50~mW cho thiết bị đeo \cite{banbury2021mlperf, ray2023tinyml_survey}. Benchmark MLPerf Tiny \cite{banbury2021mlperf} thiết lập số liệu tham chiếu: Cortex-M4F 64~MHz đạt 10,3 suy luận/giây cho phát hiện từ khóa (mô hình 49~KB), 15,2 suy luận/giây cho phát hiện bất thường (mô hình 32~KB). Các chiến lược tối ưu hóa chính:
\begin{itemize}
    \item \textbf{Lượng tử hóa}: Float 32-bit sang số nguyên 8-bit giảm kích thước 4$\times$, tăng tốc suy luận 2--3$\times$ với suy giảm 0,5--2\% \cite{lin2022mcunet}.
    \item \textbf{Cắt tỉa}: Loại bỏ 40--70\% trọng số (theo độ lớn hoặc cấu trúc) cho tăng tốc 2--3$\times$ với mất dưới 1\% \cite{schmidhuber2015deep}.
    \item \textbf{Chưng cất kiến thức}: Mô hình nhỏ (10--20\% tham số) bắt chước mô hình lớn, đạt 85--95\% hiệu suất với kích thước nhỏ hơn 5--10$\times$ \cite{dutta2023tinyml_platforms}.
    \item \textbf{ML cổ điển}: Random Forests (độ sâu $\leq$10) và SVM (100--500 support vectors) đạt 88--93\% với 50--200~KB bộ nhớ, suy luận 12--25~ms \cite{wang2024har_comparative}.
\end{itemize}
Các khảo sát TinyML \cite{ray2023tinyml_survey, dutta2023tinyml_platforms} nhận diện ba mẫu triển khai: (1) ưu tiên độ trễ dùng ML cổ điển (5--20~ms), (2) ưu tiên độ chính xác dùng CNN/LSTM nhẹ (40--80~ms, cao hơn 1--2\%), (3) ưu tiên điện năng dùng lấy mẫu thưa và lượng tử hóa mạnh. Thành phần sinh mã của framework thực hiện xuất mã C/C++ theo nền tảng với chuyển đổi tham số phù hợp ràng buộc, và các chế độ tối ưu hóa đa mục tiêu cho phép lựa chọn có căn cứ giữa độ chính xác, tốc độ và tiêu thụ điện năng.

\section{Các khoảng trống trong công trình hiện có}
Mặc dù hệ sinh thái công cụ edge AI đã khá hoàn thiện, một số khoảng trống quan trọng vẫn tồn tại:
\begin{enumerate}
    \item \textbf{Tiền xử lý tương tác hạn chế}: Ít hệ thống cho phép kiểm tra trực quan tín hiệu thô so với đã lọc và tinh chỉnh phân đoạn thủ công trong cùng môi trường.
    \item \textbf{Thiếu minh bạch trong trích xuất đặc trưng}: Các công cụ tạo đặc trưng tự động che giấu bước dẫn xuất, cản trở mục đích giáo dục và tái tạo nghiên cứu.
    \item \textbf{Sinh mã triển khai thiếu linh hoạt}: Mã đầu ra thường là khối nguyên, khó tùy chỉnh hoặc thích ứng theo phần cứng cụ thể.
    \item \textbf{Đánh đổi giữa dễ sử dụng và kiểm soát}: Các nền tảng hiện tại thường thiên về một cực---dễ dùng như Teachable Machine, hoặc linh hoạt cao như TensorFlow---ít nền tảng cân bằng được cả hai cho quy trình xử lý dữ liệu cảm biến chuyển động.
\end{enumerate}
Framework đề xuất giải quyết các khoảng trống này bằng cách thống nhất tiền xử lý tương tác, trích xuất đặc trưng minh bạch, huấn luyện modular và sinh mã tối ưu hóa dưới một kiến trúc mở rộng.

\section{Tóm tắt}
Các nghiên cứu và nền tảng công nghiệp trước đây cung cấp nền tảng vững chắc cho theo dõi chuyển động, nhưng chưa có nền tảng nào cung cấp quy trình tích hợp, minh bạch và thân thiện với người dùng, được tối ưu hóa cho triển khai biên dựa trên IMU với khả năng phân đoạn dữ liệu tương tác. Công trình này xây dựng trên các thuật toán và thực hành edge ML đã được thiết lập, đồng thời đóng góp phương pháp mới giúp phổ cập quy trình phát triển, tăng tốc lặp lại và cải thiện tính sẵn sàng triển khai cho các ứng dụng theo dõi chuyển động thực tế.
